\documentclass[10pt,a4paper]{amsart}
\usepackage[margin=19mm]{geometry}
\usepackage{amsmath,amssymb,mathtools}
\usepackage[scr=rsfs,cal=euler]{mathalfa}
\usepackage{xfrac}
\usepackage[unicode,hidelinks,bookmarks=false]{hyperref}
\newcommand{\ie}{i.e.\ }
\newcommand{\cf}{cf.\ }
\newcommand{\resp}{respectively\ }
\newcommand{\Subsection}{\subsection}
\providecommand{\nobreakdash}{\nobreak\hbox{-}}
\setlength{\emergencystretch}{2em}
\multlinegap0pt
\usepackage{amssymb}
\newtheorem{lemma}{Lemma}
\newtheorem{theorem}{Theorem}
\title{Fractional Powers of Operators: Characterization via $\delta$-Regularized Logarithmic Representation}
\author{Yoritaka Iwata}
\date{Source: arXiv:2608.13903v1 (CC BY 4.0)}
\begin{document}
\maketitle

\noindent\emph{Source and licence.} Fractional Powers of Operators: Characterization via $\delta$-Regularized Logarithmic Representation. Version arXiv:2608.13903v1 (CC BY 4.0), distributed by arXiv under the Creative Commons Attribution 4.0 International licence, http://creativecommons.org/licenses/by/4.0/, as recorded in the arXiv licence metadata for this exact version and shown on the abstract page https://arxiv.org/abs/2608.13903v1. Full licence notice: \texttt{licenses/arxiv-2608.13903-CC-BY-4.0.txt}.\par\smallskip

\noindent\textit{Editorial scope.} Counted unit: the article's theorem thm-01 (source0.tex lines 207--225). The article's complete original proof of it is delivered with it (source0.tex lines 228--265, one \texttt{proof} environment of 38 lines). No mathematics is written, repaired or re-derived: the statement and the proof are the article's own lines, in the article's own notation, and no retained line is substituted or re-flowed.  Retained uncounted as the source-local context the proof needs: lines 156--165.  Nothing was omitted from any retained span, so the package carries no omission record.  No retained line contains a citation, so the wrapper carries no bibliography and no bibliography entry.  No retained line contains a figure environment, an image-inclusion command or drawing code, so no image is delivered and the package declares no asset dependency.  Editorial formatting only, in addition: an AMS \texttt{amsart} preamble that restates the article's own declarations and macros used by the retained lines, namely the environments \texttt{lemma}, \texttt{theorem} on separate counters, together with the standard packages those lines need.  The retained environments print in this document as Lemma 1, Theorem 1 in order of appearance, whereas the article prints the same items with the numbering of its own full document.  The complete native source of the article as distributed in the arXiv e-print for this exact version is bundled unchanged as \texttt{sources/207622/source0.tex}.
\begin{lemma}[Representation of Infinitesimal Generators of One-Parameter Semigroups] \label{lem03}
Let $\{U(\lambda)\}_{\lambda \in [a, b]}$ be a one-parameter $C^0$-semigroup on a Banach space $X$ with its infinitesimal generator $D$ satisfying $0 \in \rho(D)$. 
Let $I_{\eta}(\lambda) = (I - \eta^{-1}U(\lambda))^{-1}$ be the family of resolvent operators determined by $U(\lambda)$ for a properly chosen complex parameter $\eta \in \mathbb{C}$. 
Let the operator $\mathcal{D}$ be defined by
\begin{equation} \label{eq-simp01}
\mathcal{D} := \lambda^{-1} \left( \operatorname{Log}[\eta(I_{\eta}-I)] - \operatorname{Log}(I_{\eta}) \right),
\end{equation}
where $\operatorname{Log}$ denotes the principal value of the operator logarithm. 
Then, the operator $\mathcal{D}$ acts as an infinitesimal generator explicitly independent of $\lambda$, coinciding with $D$ on its domain $D(D)$, and generates the one-parameter $C^0$-semigroup $U(\lambda)$ on $X$.
\end{lemma}

\begin{theorem} [$\delta$-Regularization of Operator] \label{thm-01}
Let $X$ be a Banach space, and let $\{U(\lambda)\}_{\lambda \ge 0}$ be a $C^0$-semigroup on $X$ generated by an infinitesimal generator $D$ whose resolvent set contains the origin ($0 \in \rho(D)$). 
For the unique solution $u(\lambda) = U(\lambda-a)u_a$ of the abstract Cauchy problem:
\begin{equation} \label{cauchy-main}
d u(\lambda) /d \lambda  = Du(\lambda), \quad \lambda \in (a, b], 
\quad  u(a) = u_a \in X. 
\end{equation}
Let $\mathcal{D}$ be the logarithmic representation of $D$ given by Lemma~\ref{lem03}:
\[
\mathcal{D} = \lambda^{-1} \left( \operatorname{Log}[\eta(I_{\eta}-I)] - \operatorname{Log}[I_{\eta}] \right),
\]
where $I_{\eta} = (I - \eta^{-1} U(\lambda))^{-1}$ for a sufficiently large parameter $\eta > 0$.

Then, for any $k \in \mathbb{R}$, the fractional power $D^k$ is well-defined via the $\delta$-regularized logarithmic representation $\mathcal{D}^k$ as
\begin{equation} \label{fractional-def}
\mathcal{D}^k = \exp \left[ k \left\{ \operatorname{Log} \left[ I - \delta (\mathcal{D} + \delta I)^{-1} \right] + \operatorname{Log} [\mathcal{D} + \delta I] \right\} \right],
\end{equation}
where $\delta > 0$ is a positive regularization parameter chosen such that the operator $(\mathcal{D} + \delta I)$ possesses a bounded inverse and the operator logarithms are uniquely defined within their principal branches.
\end{theorem}

\begin{proof}
The validity of $\delta$-regularization is established.
In particular, starting from the logarithmic representation of the bounded operator $\mathcal{D}$ constructed in Lemma~\ref{lem03}, the definition of the fractional power operator $\mathcal{D}^k$ is shown to be uniquely determined, completely independent of the choice of the regularization parameter $\delta > 0$, and well-defined on $X$.\vspace{2mm}\\
%%%%%%
\noindent\textbf{Step 1 (Algebraic Setup via $\delta$-Regularization).}
In general, the infinitesimal generator $D$ of a $C^0$-semigroup is unbounded, and domain-related analytical constraints pose severe difficulties when directly defining its operator logarithm or fractional powers.
To overcome this issue, an algebraic $\delta$-regularization with a parameter $\delta > 0$ is applied to the family of bounded operators. 
Specifically, by exploiting the formal algebraic identity
\begin{equation} \label{algebraic-decomp}
\mathcal{D} = (\mathcal{D} + \delta I) \left[ I - \delta (\mathcal{D} + \delta I)^{-1} \right],
\end{equation}
the $\delta$-regularization avoids directly manipulating $\mathcal{D}$, which may exhibit unbounded nature or spectral singularities, by reconstructing the fractional power operators with the shifted operator $\mathcal{D} + \delta I$ as the fundamental generating unit.\vspace{2mm}\\
%%%%%%
\noindent\textbf{Step 2 (Well-Definedness via Functional Calculus).} 
By assumption, the regularization parameter $\delta > 0$ is chosen such that $(\mathcal{D} + \delta I)$ possesses a bounded inverse. 
Since $\mathcal{D}$ is a bounded operator, the spectral mapping theorem implies that the spectrum $\sigma(\mathcal{D} + \delta I)$ is a compact set in the complex plane. 
Furthermore, by selecting $\delta$ appropriately, this spectrum is strictly contained in the domain of the principal branch of the logarithm ($\mathbb{C} \setminus (-\infty, 0]$). 
Consequently, by the Dunford--Riesz functional calculus for bounded operators, $\operatorname{Log} [\mathcal{D} + \delta I]$ is uniquely and well-defined as a bounded operator in $B(X)$, and its spectrum $\sigma(\operatorname{Log}[\mathcal{D} + \delta I])$ is also bounded. 
Similarly, applying the exponential mapping, the spectral mapping theorem guarantees that the operator $\exp(k \operatorname{Log} [\mathcal{D} + \delta I])$ remains well-defined and bounded in $B(X)$ for any $k \in \mathbb{R}$.
\vspace{2mm}\\
%%%%%%
\noindent\textbf{Step 3 (Independence of $\delta$ and Conclusion).} 
The first term inside the exponential function in the definition \eqref{fractional-def} is examined:
\[ A_\delta := I - \delta (\mathcal{D} + \delta I)^{-1}.\]
Due to the boundedness of $(\mathcal{D} + \delta I)^{-1}$, $A_\delta$ is a bounded operator on $X$, and simple algebraic manipulation yields $A_\delta = \mathcal{D}(\mathcal{D} + \delta I)^{-1}$. 
According to the spectral mapping theorem, its spectrum is given by
\[
\sigma(A_\delta) = \left\{ \frac{z}{z + \delta} \;\middle|\; z \in \sigma(\mathcal{D}) \right\}.
\]
Since $0$ belongs to the resolvent set $\rho(\mathcal{D})$ by assumption, it follows that $0 \in \rho(A_\delta)$. 
This guarantees that $\sigma(A_\delta)$ is bounded away from the origin and lies strictly within the domain of the principal branch of the logarithm. 
Consequently, by the logarithmic property for commuting operators, the operational logarithm in Eq.~\eqref{fractional-def} simplifies to
\[ 
k \left\{ \operatorname{Log} \left[ \mathcal{D}(\mathcal{D} + \delta I)^{-1} \right] + \operatorname{Log} [\mathcal{D} + \delta I] \right\} = k \operatorname{Log} \mathcal{D}. 
\]
This demonstrates that the fractional power operator $\mathcal{D}^k$ is determined solely by the intrinsic properties of $\mathcal{D}$, completely independent of the choice of the regularization parameter $\delta > 0$. \vspace{2mm} \\
As a result, the fractional power operator $\mathcal{D}^k$ is proven to be uniquely well-defined as a bounded operator on $X$, and can be consistently applied to the solution $u(\lambda)$ of the abstract Cauchy problem.\vspace{2mm}\\
\end{proof}

\end{document}
